\documentclass[11pt,a4paper]{article}
\usepackage{_preamble}

\title[Conditions Générales applicables aux Balades \\ version du \today]{Conditions Générales applicables aux Balades}
\subtitle{Annexe I au \textit{Règlement Intérieur} (\no 25--02/CES/DIR)}
\author{}
\date{\today}

\begin{document}
   \annexHeader{}
   \maketitle

   \vspace{-30pt}

   \subsection*{Balades à poney}
   En juillet et août, les balades à poney se font sur réservation, à horaires définis.
   De durée fixe, elles se font sous encadrement du Centre Équestre, le poney étant tenu par un parent ou responsable majeur.
   Le reste du temps, les tours en main se font sous l'entière responsabilité des parents ou leur délégué majeur.

   Dans tous les cas, elles se font \emph{uniquement au pas}, à raison d'un enfant par poney, équipé d'un casque homologué.

   L'accompagnant familial doit disposer de toutes ses facultés physiques pour assurer la sécurité de l'enfant à poney lors de la prise en main, au risque de se voir refuser la balade.
   Il est demandé aux familles de prévoir \emph{1 adulte par poney}.
   Un départ en balade peut être refusé si toutes les conditions permettant son déroulement en toute sécurité ne sont pas réunies.

   \subsection*{Balades à cheval}
   Les balades à cheval se font seulement sur réservation.
   Il peut être demandé, pour valider l'inscription, un paiement total ou partiel de la prestation.

   Toute balade à cheval est encadrée par le nombre d'enseignants défini par la direction, en accord avec le cadre réglementaire en vigueur.
   Les enseignants sont responsables de leur groupe, et sont habilités à prendre toute décision nécessaire au maintien de la sécurité aussi bien des cavaliers que du groupe dans son ensemble.
   \\

   On distingue, en balade, deux groupes distincts : les \og débutants \fg{}, toute personne non expérimentée montant à cheval pour la première fois ou seulement occasionnellement, ou de niveau incompatible avec le galop, et les \og confirmés \fg{} qui sont les cavaliers possédant un niveau équestre (du point de vue fédéral) et étant en capacités de trotter et galoper en toute aisance en extérieur\footnote{Galop 3 ou équivalent.} sans atteinte à la sécurité.
   Tout cavalier ne répondant pas à ces critères est \textit{de facto} considéré comme débutant.

   Le groupe constitué des cavaliers débutants fera la balade \emph{au pas uniquement}.
   Au contraire une balade de niveau confirmé comportera des phases de trot et de galop, \emph{de longueur et durée décidées par l'enseignant qui est le seul à avoir connaissance des capacités et de l'état de sa cavalerie} et du niveau des cavaliers qu'il encadre.

   Les enseignants ayant pour mission première de garantir la sécurité du groupe.
   À ce titre, ils ont la possibilité de refuser des cavaliers dans le groupe confirmé s'ils estiment que trotter ou galoper présente un danger.
   \emph{La responsabilité des enseignants ne peut être engagée en cas de chute d'un cavalier n'ayant pas respecté cette décision, ou toute autre consigne de sécurité donnée avant ou pendant la balade}.

   \subsection*{Tenue exigée et aptitude physique}
   Aucune tenue spécifique n'est exigée, cependant il est conseillé aux cavaliers, pour leur sécurité et leur confort, de porter un pantalon long et des chaussures fermées.
   Néanmoins, pour des raisons de sécurité, l'enseignant se réserve le droit de refuser un cavalier portant une tenue non adaptée à la pratique de l'équitation en extérieur (tongs, sac à dos, écharpe, déguisement\dots).

   Quiconque monte en balade, quel que soit son niveau, est pleinement conscient qu'\emph{il s'agit d'une pratique sportive d'extérieur}.
   Aussi, il est de la responsabilité de chacun de s'assurer qu'il est en pleine capacité de pratiquer l'équitation, même au pas et en balade.

   En particulier, \emph{toute condition médicale ou physique incompatible avec l'activité proposée est supposée connue du pratiquant}, qui monte alors en pleine connaissance des risques éventuels du fait de sa condition.
   En cas de doute sur la pleine aptitude à pratiquer l'équitation en extérieur, les enseignants peuvent refuser le départ d'un cavalier, pour des raisons de sécurité individuelle et collective.

   Même si nous faisons tout pour le réduire, \emph{le risque de chute existe}.

   \subsection*{Handicap}
   Nous défendons l'équitation pour tous, et souhaitons proposer à chacun une activité adaptée à ses capacités.
   À ce titre, il découle du bon sens que \emph{tout handicap ou situation nécessitant un accompagnement particulier doit être signalé au moment de la réservation}, afin que nous puissions nous organiser au mieux pour vous accueillir et adapter la cavalerie.

   Dans le cas contraire, la direction peut refuser le départ si elle estime que les moyens humains ou techniques sont insuffisants pour garantir la sécurité.

   \subsection*{Annulation et remboursement de la prestation}
   De manière générale, un faux-départ n'est pas remboursé.

   Les avances perçues au titre d'une réservation de balade à cheval sont systématiquement rendues en cas d'annulation, \emph{au plus tard deux heures avant l'heure de départ}.
   Passé ce délai, elles sont dues (hors cas de force majeure).
   Elles sont systématiquement rendues en cas d'annulation par le Centre Équestre.\\

   Aucun remboursement n'est possible pour un cavalier auquel aurait été refusé le galop.

   Les cartes cadeau, quelles que soient leur valeur et leur validité, ne pourront pas être remboursées ni échangées.
   De plus, leur validité ne peut être repoussée hors cas de force majeure.

   \vfill

   \hspace{.55\textwidth}
   \begin{minipage}{.4\textwidth}
      \begin{center}
         Stéphanie \textsc{Marais},\\
         Dirigeant le Centre Équestre de Sauveterre
      \end{center}
   \end{minipage}

\end{document}
